\documentclass[10pt, aspectratio=169]{beamer}
% Preámbulo modular para presentaciones Beamer (Modelación Estadística)
% Diseño y paleta institucional del Tecnológico de Monterrey

\documentclass[aspectratio=169,xcolor={dvipsnames,table}]{beamer}

\usepackage[utf8]{inputenc}
\usepackage[spanish,mexico]{babel}
\usepackage{amsmath,amssymb,amsfonts,mathtools}
\usepackage{graphicx}
\usepackage{listings}
\usepackage{booktabs}
\usepackage{tikz}
\usetikzlibrary{matrix,arrows,decorations.pathmorphing,shapes.geometric,positioning}

% Paleta Institucional Tecnológico de Monterrey
\definecolor{TecAzul}{HTML}{0039A6}       % Azul TEC: color primario institucional
\definecolor{TecAzulOscuro}{HTML}{1A2E51} % Azul marino oscuro
\definecolor{TecRojo}{HTML}{EC2661}       % Rojo de acento (paleta miTec)
\definecolor{TecGris}{HTML}{646464}       % Gris neutro para comentarios y subtítulos
\definecolor{TecFondoBloque}{HTML}{F4F6F9}% Fondo suave para bloques

% Configuración de Tema Beamer
\usetheme{Boadilla}
\usecolortheme[named=TecAzulOscuro]{structure}
\setbeamercolor{alerted text}{fg=TecRojo}
\setbeamercolor{palette primary}{bg=TecAzulOscuro,fg=white}
\setbeamercolor{palette secondary}{bg=TecAzul,fg=white}
\setbeamercolor{palette tertiary}{bg=TecAzulOscuro!80!black,fg=white}
\setbeamercolor{title}{fg=TecAzulOscuro,bg=TecFondoBloque}
\setbeamercolor{frametitle}{fg=TecAzulOscuro,bg=TecFondoBloque}

% Bloques personalizados elegantes
\setbeamercolor{block title}{bg=TecAzulOscuro,fg=white}
\setbeamercolor{block body}{bg=TecFondoBloque,fg=black}
\setbeamercolor{block title alerted}{bg=TecRojo,fg=white}
\setbeamercolor{block body alerted}{bg=TecRojo!8,fg=black}
\setbeamercolor{block title example}{bg=TecAzul,fg=white}
\setbeamercolor{block body example}{bg=TecAzul!8,fg=black}

% Redondear bordes y añadir sombra sutil
\setbeamertemplate{blocks}[rounded][shadow=true]
\setbeamertemplate{navigation symbols}{}

% Pie de página limpio con número de diapositiva
\setbeamertemplate{footline}{
  \leavevmode%
  \hbox{%
  \begin{beamercolorbox}[wd=.7\paperwidth,ht=2.25ex,dp=1ex,left]{author in head/foot}%
    \hspace*{2ex}\insertshortauthor~~(\insertshortinstitute) \hfill \insertshorttitle
  \end{beamercolorbox}%
  \begin{beamercolorbox}[wd=.3\paperwidth,ht=2.25ex,dp=1ex,right]{date in head/foot}%
    \insertshortdate{}\hspace*{2em}
    \insertframenumber{} / \inserttotalframenumber\hspace*{2ex}
  \end{beamercolorbox}}%
  \vskip0pt%
}

% Configuración de Listings de Python para visualización pedagógica de código
\lstset{
    language=Python,
    basicstyle=\ttfamily\footnotesize,
    keywordstyle=\color{TecAzulOscuro}\bfseries,
    stringstyle=\color{TecRojo},
    commentstyle=\color{TecGris}\itshape,
    showstringspaces=false,
    breaklines=true,
    frame=lines,
    rulecolor=\color{TecAzulOscuro!30},
    backgroundcolor=\color{TecFondoBloque},
    aboveskip=0.5em,
    belowskip=0.5em,
    literate={á}{{\'a}}1 {é}{{\'e}}1 {í}{{\'i}}1 {ó}{{\'o}}1 {ú}{{\'u}}1
             {Á}{{\'A}}1 {É}{{\'E}}1 {Í}{{\'I}}1 {Ó}{{\'O}}1 {Ú}{{\'U}}1
             {ñ}{{\~n}}1 {Ñ}{{\~N}}1 {¿}{{?`}}1 {¡}{{!`}}1
}

% Metadatos institucionales por defecto
\title[Modelación Estadística]{Modelación Estadística para Ciencia de Datos e Ingeniería}
\author[J. Castillo Colmenares]{Juliho David Castillo Colmenares}
\institute[Tec de Monterrey]{Tecnológico de Monterrey \\ Escuela de Ingeniería y Ciencias}
\date{\today}

% Comandos y macros estadísticas compartidas para las presentaciones Beamer (Metropolis)

% Conjuntos numéricos básicos
\newcommand{\R}{\mathbb{R}}
\newcommand{\N}{\mathbb{N}}
\newcommand{\Q}{\mathbb{Q}}
\newcommand{\Z}{\mathbb{Z}}
\newcommand{\set}[1]{\left\{ #1 \right\}}
\newcommand{\abs}[1]{\left|#1\right|}
\newcommand{\norm}[1]{\left\| #1 \right\|}

% Comandos estadísticos y probabilísticos del curso
\newcommand{\Var}{\operatorname{Var}}
\newcommand{\cov}{\operatorname{Cov}}
\newcommand{\corr}{\rho}
\newcommand{\comb}[2]{\begin{pmatrix} #1 \\ #2 \end{pmatrix}}
\newcommand{\s}{\sigma}
\newcommand{\particion}{\mathfrak{P}}
\newcommand{\card}[1]{\#\left( #1 \right)}
\newcommand{\seq}[3]{\set{#1}_{#2}^{#3}}
\newcommand{\E}{\mathbb{E}}
\newcommand{\Prob}{\mathbb{P}}

% Entornos pedagógicos y cajas en español académico natural
\newenvironment{discusion}{%
  \begin{alertblock}{Pregunta de Discusión}%
}{%
  \end{alertblock}%
}

\newenvironment{reflexion}{%
  \begin{alertblock}{Reflexión para el Aula}%
}{%
  \end{alertblock}%
}

\newenvironment{paradoja}{%
  \begin{alertblock}{Paradoja de Exploración}%
}{%
  \end{alertblock}%
}

\newenvironment{motivacion}{%
  \begin{exampleblock}{Motivación: Ciencia de Datos en la Práctica}%
}{%
  \end{exampleblock}%
}

\newenvironment{retoclase}{%
  \begin{block}{Reto Rápido en Equipo}%
}{%
  \end{block}%
}


\title[Gamma, Beta y Weibull]{Sección 04.07: Distribuciones Gamma, Beta y Weibull}
\subtitle{Familia Gamma, Proporciones en $[0,1]$ y Confiabilidad}
\author[J. Castillo Colmenares]{Juliho Castillo Colmenares}
\institute{Tecnológico de Monterrey}
\date{\vspace{-1.2cm}}

\begin{document}

% Slide 1: Portada
\begin{frame}[plain]
  \vspace{-0.8cm}
  \titlepage
\end{frame}

% Slide 2: Hoja de Ruta
\begin{frame}{Hoja de Ruta --- Capítulo 04: Variables Aleatorias Continuas}
  \footnotesize
  ¿Dónde nos encontramos en el desarrollo modular del capítulo? \pause
  \begin{itemize}\itemsep=0.08em
    \item \textbf{04.01} Función de Densidad (PDF) y Soporte Continuo \pause
    \item \textbf{04.02} Función de Distribución Acumulada (CDF) Continua y Cuantiles \pause
    \item \textbf{04.03} Esperanza, Varianza y LOTUS Continuo \pause
    \item \textbf{04.04} Distribución Uniforme Continua $U(a, b)$ \pause
    \item \textbf{04.05} Distribución Normal / Gaussiana y Puntaje $Z$ \pause
    \item \textbf{04.06} Distribución Exponencial y Procesos sin Memoria \pause
    \item \textbf{04.07} Distribuciones Gamma, Beta y Weibull \emph{(Hoy --- Cierre del Capítulo)}
  \end{itemize}
  \vspace{-0.05cm}
  \begin{block}{Objetivo de la Sesión}
    Generalizar la familia Gamma más allá de la Exponencial y la Chi-cuadrada, introducir la distribución Beta para modelar proporciones en $[0,1]$, y dominar la distribución Weibull como herramienta central del análisis de confiabilidad y tasas de falla variables.
  \end{block}
\end{frame}

% Slide 3: Motivación
\begin{frame}{Tres Formas Flexibles Más Allá de lo Simétrico}
  \footnotesize
  \begin{motivacion}
    Hasta ahora hemos visto distribuciones con formas fijas: uniforme (plana), exponencial (decreciente) y normal (simétrica). Muchos fenómenos reales requieren formas más flexibles: tiempos de espera hasta el $\alpha$-ésimo evento, proporciones acotadas en $[0,1]$, y tiempos de falla cuyo riesgo cambia con la edad del componente.
  \end{motivacion}

  \vspace{0.15cm}
  \pause
  \begin{itemize}\itemsep=0.1cm
    \item \textbf{Gamma:} Generaliza la Exponencial permitiendo sumar $\alpha$ tiempos de espera; contiene a la Chi-cuadrada como caso particular. \pause
    \item \textbf{Beta:} La distribución canónica en $[0,1]$; prior conjugado natural para proporciones bayesianas. \pause
    \item \textbf{Weibull:} Generaliza la Exponencial permitiendo tasas de falla crecientes o decrecientes; pilar de la ingeniería de confiabilidad.
  \end{itemize}
\end{frame}

% Slide 4: Definición Formal - Gamma
\begin{frame}{Definición Formal de la Distribución Gamma}
  \footnotesize
  Una variable aleatoria continua $X$ tiene distribución \emph{Gamma} con parámetros de forma $\alpha > 0$ y escala $\beta > 0$ si su PDF es: \pause

  \vspace{0.1cm}
  \begin{block}{Función de Densidad de Probabilidad}
    $$f(x) = \dfrac{1}{\Gamma(\alpha)\,\beta^{\alpha}}\,x^{\alpha-1}\,e^{-x/\beta}, \quad x > 0$$
    donde $\Gamma(\alpha) = \int_{0}^{\infty} t^{\alpha-1}e^{-t}\,dt$ es la función gamma. Escribimos $X \sim \text{Gamma}(\alpha, \beta)$.
  \end{block}

  \vspace{0.1cm}
  \pause
  \begin{alertblock}{Propiedades Fundamentales}
    \begin{itemize}\itemsep=0.05cm
      \item Media y varianza: $\mu_{X} = \alpha\beta$, $\sigma_{X}^{2} = \alpha\beta^{2}$.
      \item MGF: $M_{X}(t) = (1-\beta t)^{-\alpha}$, para $t < 1/\beta$.
      \item \textbf{Propiedad aditiva:} si $X_{i} \sim \text{Gamma}(\alpha_{i}, \beta)$ independientes (misma escala), $\sum X_{i} \sim \text{Gamma}(\sum \alpha_{i}, \beta)$.
    \end{itemize}
  \end{alertblock}
\end{frame}

% Slide 5: Casos particulares
\begin{frame}{Casos Particulares: Exponencial y Chi-cuadrada}
  \footnotesize
  La familia Gamma unifica varias distribuciones ya conocidas: \pause

  \vspace{0.1cm}
  \begin{block}{Exponencial y Chi-cuadrada como Casos Gamma}
    \begin{itemize}\itemsep=0.05cm
      \item $\text{Gamma}(1, \beta) = \text{Exp}(\lambda = 1/\beta)$: espera hasta el \emph{primer} evento de un proceso Poisson.
      \item $\text{Gamma}(\nu/2, 2) = \chi^{2}_{\nu}$: chi-cuadrada con $\nu$ grados de libertad.
    \end{itemize}
  \end{block}

  \vspace{0.1cm}
  \pause
  \begin{alertblock}{Erlang: Suma de Tiempos de Espera}
    \begin{itemize}\itemsep=0.05cm
      \item Si $T_{1}, \dots, T_{n} \sim \text{Exp}(\lambda)$ i.i.d., el tiempo hasta el $n$-ésimo evento es $T = \sum T_{i} \sim \text{Gamma}(n, 1/\lambda)$ (distribución Erlang).
      \item Ejemplo: llamadas a un \emph{call center} con tasa $\lambda=3$/min; el tiempo hasta la tercera llamada es $\text{Gamma}(3, 1/3)$, con $P(T>1.5) \approx 0.174$.
    \end{itemize}
  \end{alertblock}
\end{frame}

% Slide 6: Distribución Beta
\begin{frame}{Definición Formal de la Distribución Beta}
  \footnotesize
  Una variable aleatoria continua $X$ tiene distribución \emph{Beta} con parámetros $\alpha, \beta > 0$ si su PDF es: \pause

  \vspace{0.1cm}
  \begin{block}{Función de Densidad de Probabilidad}
    $$f(x) = \dfrac{1}{B(\alpha,\beta)}\,x^{\alpha-1}(1-x)^{\beta-1}, \quad 0 < x < 1$$
    donde $B(\alpha,\beta) = \int_{0}^{1} t^{\alpha-1}(1-t)^{\beta-1}\,dt = \dfrac{\Gamma(\alpha)\Gamma(\beta)}{\Gamma(\alpha+\beta)}$. Escribimos $X \sim \text{Beta}(\alpha,\beta)$.
  \end{block}

  \vspace{0.1cm}
  \pause
  \begin{alertblock}{Propiedades Fundamentales}
    \begin{itemize}\itemsep=0.05cm
      \item Media y varianza: $\mu_{X} = \dfrac{\alpha}{\alpha+\beta}$, $\sigma_{X}^{2} = \dfrac{\alpha\beta}{(\alpha+\beta)^{2}(\alpha+\beta+1)}$.
      \item Caso particular: $\text{Beta}(1,1) = U(0,1)$.
      \item Simetría: $\text{Beta}(\alpha,\beta) = 1 - \text{Beta}(\beta,\alpha)$.
    \end{itemize}
  \end{alertblock}
\end{frame}

% Slide 7: Aplicación Beta
\begin{frame}{Beta como Modelo de Proporciones y Prior Bayesiano}
  \footnotesize
  La distribución Beta es la herramienta canónica para modelar variables acotadas en $[0,1]$: \pause

  \vspace{0.1cm}
  \begin{block}{Modelado de Proporciones}
    Si $X$ representa una proporción desconocida (clientes satisfechos, tasa de defectos, probabilidad de éxito), $X \sim \text{Beta}(\alpha,\beta)$ captura tanto el valor esperado como la incertidumbre alrededor de él.
  \end{block}

  \vspace{0.1cm}
  \pause
  \begin{alertblock}{Prior Conjugado Bayesiano}
    \begin{itemize}\itemsep=0.05cm
      \item Si el prior sobre $p$ (parámetro Bernoulli/Binomial) es $\text{Beta}(\alpha,\beta)$ y se observan $k$ éxitos en $n$ ensayos, el posterior es $\text{Beta}(\alpha+k,\, \beta+n-k)$.
      \item Ejemplo: prior $\text{Beta}(2,2)$, $n=20$, $k=14$ éxitos $\implies$ posterior $\text{Beta}(16,8)$, con media posterior $2/3$.
    \end{itemize}
  \end{alertblock}
\end{frame}

% Slide 8: Distribución Weibull
\begin{frame}{Definición Formal de la Distribución Weibull}
  \footnotesize
  Una variable aleatoria $T \geq 0$ (tiempo de vida) tiene distribución \emph{Weibull} con forma $\beta > 0$ y escala $\eta > 0$ si su PDF es: \pause

  \vspace{0.1cm}
  \begin{block}{Función de Densidad de Probabilidad}
    $$f(t) = \dfrac{\beta}{\eta}\left(\dfrac{t}{\eta}\right)^{\beta-1}\exp\!\left[-\left(\dfrac{t}{\eta}\right)^{\beta}\right], \quad t \geq 0$$
    con $F(t) = 1-\exp[-(t/\eta)^{\beta}]$. Escribimos $T \sim \text{Weibull}(\beta,\eta)$.
  \end{block}

  \vspace{0.1cm}
  \pause
  \begin{alertblock}{Tasa de Falla $h(t) = \dfrac{\beta}{\eta}(t/\eta)^{\beta-1}$}
    \begin{itemize}\itemsep=0.05cm
      \item $\beta = 1$: tasa constante, $T \sim \text{Exp}(1/\eta)$ (sin envejecimiento).
      \item $\beta < 1$: tasa decreciente (mortalidad infantil, defectos de fabricación).
      \item $\beta > 1$: tasa creciente (desgaste mecánico acumulado).
    \end{itemize}
  \end{alertblock}
\end{frame}

% Slide 9: Aplicaciones
\begin{frame}{Aplicaciones en Ingeniería, Confiabilidad y Ciencia de Datos}
  \footnotesize
  Las tres distribuciones son pilares de la modelación aplicada: \pause

  \vspace{0.15cm}
  \begin{itemize}\itemsep=0.12cm
    \item \textbf{Confiabilidad Industrial:} Weibull modela tiempos de falla en manufactura, mantenimiento predictivo y garantías de producto. \pause
    \item \textbf{Teoría de Colas:} Gamma/Erlang modela tiempos de servicio con múltiples etapas ($M/E_{k}/1$). \pause
    \item \textbf{Estadística Bayesiana:} Beta es el prior conjugado estándar para tasas de conversión, tasas de error y proporciones en A/B testing. \pause
    \item \textbf{Ciencia de Datos:} Regularización Bayesiana, modelos jerárquicos y análisis de supervivencia (\emph{survival analysis}) en salud e ingeniería.
  \end{itemize}
\end{frame}

% Slide 10: Lab Python (Bloque 1)
\begin{frame}[fragile]{Laboratorio en Python: Gamma y Propiedad Erlang}
  \scriptsize
  Validación de $\int f = 1$ para la Gamma, propiedad aditiva de Erlang, y casos particulares Exponencial/Chi-cuadrada (\texttt{04.07\_gamma\_beta\_weibull.py}):
  {\lstset{basicstyle=\fontsize{3.6pt}{4.3pt}\selectfont\ttfamily,aboveskip=0.1em,belowskip=0.1em}
  \lstinputlisting[firstline=17, lastline=51]{../../code/04_variables_aleatorias_continuas/04.07_gamma_beta_weibull.py}}
\end{frame}

% Slide 11: Lab Python (Bloque 2)
\begin{frame}[fragile]{Laboratorio en Python: Beta, Momentos y Actualización Bayesiana}
  \scriptsize
  Cálculo de momentos, verificación de la propiedad de simetría, y actualización de un prior conjugado:
  \vspace{0.02cm}
  {\lstset{basicstyle=\fontsize{4pt}{4.8pt}\selectfont\ttfamily}
  \lstinputlisting[firstline=54, lastline=83]{../../code/04_variables_aleatorias_continuas/04.07_gamma_beta_weibull.py}}
\end{frame}

% Slide 12: Lab Python (Bloque 3)
\begin{frame}[fragile]{Laboratorio en Python: Weibull y Análisis de Confiabilidad}
  \scriptsize
  Comparación de funciones de riesgo y curvas de confiabilidad para distintos parámetros de forma:
  \vspace{0.02cm}
  {\lstset{basicstyle=\fontsize{4pt}{4.8pt}\selectfont\ttfamily}
  \lstinputlisting[firstline=86, lastline=115]{../../code/04_variables_aleatorias_continuas/04.07_gamma_beta_weibull.py}}
\end{frame}

% Slide 13: Salida Terminal
\begin{frame}[fragile]{Laboratorio en Python: Salida en Terminal}
  \scriptsize
  Salida estándar generada al ejecutar el script del laboratorio en Python:
  \vspace{0.02cm}
  \begin{block}{Salida Estándar en Consola}
    \fontsize{4pt}{5pt}\selectfont
    \begin{verbatim}
=== Block 1: Gamma Distribution & Erlang ===
Gamma(alpha=3, beta=2) PDF integral: 1.00000000
  Mean: 6.0000, Variance: 12.0000
Erlang: Gamma(3,0.3333); P(T>1.5)=0.173578
Exponential case: Gamma(1,4) vs Exp(0.25): P(X>6)=0.223130 (match)
Chi-squared case: nu=3, P(X>7.815)=0.049994 (match)

=== Block 2: Beta Distribution & Moments ===
Beta(alpha=3, beta=5) PDF integral: 1.00000000
  Mean: 0.375000, Variance: 0.026042
Symmetry: P(X<=0.3)|Beta(2,5) = 1-P(Y<=0.7)|Beta(5,2) = 0.579825
Bayesian update: prior Beta(2,2) + 14/20 -> Beta(16,8), mean=0.666667

=== Block 3: Weibull & Reliability Analysis ===
Weibull(shape=1.5, scale=4) PDF integral: 1.00000000
  R(3) = 0.522297
Hazard: Lot A (beta=1) h=0.100 constant; Lot B (beta=2) h(5)=0.100, h(15)=0.300
    \end{verbatim}
  \end{block}
\end{frame}

% Slide 18: Cierre
\begin{frame}{Cierre de Sección: Gamma, Beta y Weibull}
  \begin{reflexion}
    Las distribuciones Gamma, Beta y Weibull completan nuestro estudio de las variables aleatorias continuas fundamentales. La Gamma generaliza los tiempos de espera, la Beta modela proporciones en $[0,1]$, y la Weibull captura tasas de falla variables en el tiempo. Juntas cubren la gran mayoría de los fenómenos de espera, proporción y confiabilidad en la práctica.
  \end{reflexion}

  \vspace{0.2cm}
  \pause
  \begin{block}{Lecciones Clave}
    \begin{itemize}\itemsep=0.08cm
      \item \textbf{Gamma:} $f(x) \propto x^{\alpha-1}e^{-x/\beta}$; generaliza Exponencial y Chi-cuadrada; aditiva bajo suma.
      \item \textbf{Beta:} $f(x) \propto x^{\alpha-1}(1-x)^{\beta-1}$ en $[0,1]$; prior conjugado bayesiano estándar.
      \item \textbf{Weibull:} tasa de falla $h(t) = (\beta/\eta)(t/\eta)^{\beta-1}$; $\beta=1$ sin envejecimiento, $\beta>1$ desgaste.
    \end{itemize}
  \end{block}
\end{frame}

% Slide 19: Perspectiva Modular
\begin{frame}{Perspectiva Modular --- Cierre del Capítulo 04}
  \scriptsize
  \begin{retoclase}
    Con la Sección 04.07 concluimos el Capítulo 04: Variables Aleatorias Continuas, habiendo cubierto la PDF, CDF, esperanza/varianza, y las familias Uniforme, Exponencial, Normal, Gamma, Beta y Weibull. El Capítulo 05 --- Distribuciones Muestrales --- estudiará cómo se comportan los estadísticos (media muestral, varianza muestral) al repetir el muestreo, sentando las bases de la inferencia estadística.
  \end{retoclase}

  \vspace{0.08cm}
  \pause
  \begin{alertblock}{Preparación Recomendada}
    \begin{itemize}\itemsep=0.06cm
      \item Repasa el Teorema del Límite Central (Sección 04.05): es la base teórica de las distribuciones muestrales.
      \item Reflexiona sobre por qué $(n-1)S^{2}/\sigma^{2} \sim \chi^{2}_{n-1}$: la conexión Gamma--Chi-cuadrada de hoy reaparecerá en el Capítulo 05.
    \end{itemize}
  \end{alertblock}
\end{frame}

\end{document}
